\section{Introduction}

The standard practice in language model pre-training treats data curation as a model-agnostic step: a team chooses a perplexity filtering threshold, a deduplication aggressiveness level, and applies this recipe uniformly across all models in a training cascade. Our experiments challenge this assumption. When we filter the same FineWeb corpus at perplexity threshold $\tau{=}20$, a 14M-parameter Pythia model achieves its best HellaSwag score. When we apply the same recipe to a 31M-parameter model, it is the \emph{worst} configuration we tested --- the 31M model peaks instead at $\tau{=}50$. The optimal filtering threshold spans the full range of our experimental design space across a 2.2$\times$ scale difference.

This observation has a concrete practical implication. A perplexity threshold of $\tau{=}20$ retains only 176 of 5,000 sampled FineWeb documents (3.5\%) --- it applies extreme selectivity, keeping only the most ``standard-English,'' lowest-perplexity text. A threshold of $\tau{=}50$ retains 2,074 documents (41.5\%) --- it accepts more diverse, higher-entropy content. Practitioners who apply $\tau{=}50$ uniformly to a training cascade that includes a 14M-parameter model are training that model on 12$\times$ more data than optimal, at the cost of reduced data quality per token.

\subsection{The Problem: Scale-Independent Curation Recipes}

Data curation for language model pre-training has attracted substantial recent attention. Methods including perplexity-based filtering~\cite{Zhou2024ProX,Peng2025DataMan}, near-duplicate removal~\cite{He2024SoftDedup,Penedo2025FineWeb2}, and domain mixing~\cite{Wettig2025Organize} each demonstrate improvements in downstream benchmark performance. Yet a consistent pattern runs through this literature: each method is evaluated at a single model scale, and the optimal recipe is reported as a universal recommendation.

The deeper problem is that this implicit scale-independence assumption has never been tested in a controlled multi-scale experiment. The capacity-quality trade-off hypothesis --- grounded in Bayesian in-context learning theory~\cite{Xie2021Explanation} and capacity-constrained learning dynamics --- predicts that smaller models, unable to efficiently extract signal from high-diversity, noisy distributions, should benefit from stricter quality filtering, while larger models should leverage the diversity preserved by looser filtering. This prediction cannot be evaluated from single-scale ablations.

The gap is methodological: testing the Scale $\times$ Curation interaction requires a factorial experiment varying both model scale and curation aggressiveness simultaneously. Such experiments are computationally expensive, and prior work on scalable ablation methodology~\cite{Na2024Scalable} proposes small proxy models ($\ll$100M parameters) as a shortcut --- an approach our results show is insufficient for scale-dependent interaction studies specifically.

\subsection{Key Insight}

The optimal perplexity filtering threshold is model-scale-dependent. Small models (14M parameters) learn best from highly curated, low-diversity data; larger models (31M parameters) learn best from more diverse, less aggressively filtered data. This means that a curation recipe optimized at one scale is suboptimal --- and potentially harmful --- at another scale.

Building on this insight, we run a controlled factorial experiment directly testing the Scale $\times$ Curation interaction in language model pre-training, using real FineWeb data, real GPT-2 perplexity scoring, and standard lm-evaluation-harness evaluation.

\subsection{Contributions}

Our work makes three contributions:

\begin{enumerate}
\item \textbf{Scale-dependent optimal PPL threshold confirmed at PoC scale.} We demonstrate that $\tau^*(14\text{M}){=}20$ and $\tau^*(31\text{M}){=}50$ in a 24-run factorial experiment (3 PPL thresholds $\times$ 2 dedup levels $\times$ 2 scales $\times$ 2 seeds) on FineWeb data with Pythia-architecture models. The interaction direction is consistent across all experimental conditions.

\item \textbf{Provisional feasibility boundary for proxy models.} We show empirically that proxy models below approximately 14M parameters with a 2.2$\times$ scale ratio are insufficient to exhibit the Scale $\times$ Curation interaction in our experiments --- an important methodological observation for the scalable ablation literature.

\item \textbf{MMLU floor documented at sub-100M scale.} We confirm that MMLU 4-shot accuracy is at the random baseline (0.25) for all sub-100M models at short pre-training runs, establishing HellaSwag 0-shot as the appropriate metric for pre-training ablations at this scale.
\end{enumerate}

We organize the paper as follows: Section~\ref{sec:related} surveys related work on data curation and positions our contribution; Section~\ref{sec:methodology} describes our experimental methodology; Section~\ref{sec:setup} presents experimental setup; Section~\ref{sec:results} reports results; Section~\ref{sec:discussion} discusses implications and limitations; Section~\ref{sec:conclusion} concludes.
